% ============================================================
%  Chapter 2 --- System Architecture
% ============================================================
\chapter{System Architecture}
\label{ch:architecture}

\section{High-Level Overview}

\shiftdb{} is organized into clearly separated layers, following the classical DBMS architecture described by Hellerstein et al.~\cite{hellerstein2007architecture}.  Each layer communicates with its neighbours through narrow, well-defined interfaces.

\begin{figure}[H]
\centering
\begin{tikzpicture}[node distance=1.4cm and 2cm]
  % Frontend
  \node[component, minimum width=10cm, fill=shiftPrimary!12] (fe)
    {React Frontend (Monaco Editor, Data Grid, ER Diagram, Query Builder)};

  % REST API
  \node[component, minimum width=10cm, below=of fe] (api)
    {HTTP Server / REST API Layer (\texttt{server.h})};

  % SQL Engine
  \node[component, minimum width=4.5cm, below left=1.4cm and 0.3cm of api] (parser)
    {SQL Parser\\(\texttt{tokenizer.h}, \texttt{parser.h})};
  \node[component, minimum width=4.5cm, below right=1.4cm and 0.3cm of api] (exec)
    {Query Executor\\(\texttt{executor.h})};

  % Catalog
  \node[component, minimum width=10cm, below=3.6cm of api] (catalog)
    {Catalog Manager (\texttt{catalog.h})};

  % Storage
  \node[component, minimum width=4.5cm, below left=1.4cm and 0.3cm of catalog] (bp)
    {Buffer Pool\\(\texttt{buffer\_pool.h})};
  \node[component, minimum width=4.5cm, below right=1.4cm and 0.3cm of catalog] (table)
    {Table Manager\\(\texttt{table.h})};

  % Disk
  \node[storage, below=3.6cm of catalog, minimum width=6cm] (disk)
    {Disk Manager (\texttt{disk\_manager.h})\\[2pt]\texttt{.nxdb} files + \texttt{catalog.json}};

  % Arrows
  \draw[arr] (fe) -- (api);
  \draw[arr] (api) -- (parser);
  \draw[arr] (api) -- (exec);
  \draw[arr] (parser) -- (catalog);
  \draw[arr] (exec) -- (catalog);
  \draw[arr] (catalog) -- (bp);
  \draw[arr] (catalog) -- (table);
  \draw[arr] (bp) -- (disk);
  \draw[arr] (table) -- (disk);
  \draw[dasharr] (exec) -- (table);
\end{tikzpicture}
\caption{Layered architecture of \shiftdb{}.}
\label{fig:arch-layers}
\end{figure}

\section{Layer-by-Layer Walkthrough}

\subsection{Presentation Layer --- React Frontend}

The browser-based client communicates exclusively via HTTP/JSON with the backend REST API.  Key UI components:

\begin{table}[H]
\centering
\caption{Frontend component mapping.}
\label{tab:fe-components}
\begin{tabular}{l l p{8cm}}
\toprule
\textbf{Component} & \textbf{File} & \textbf{Responsibility} \\
\midrule
SQL Editor       & \code{SqlEditor.tsx}            & Monaco-based editor with IntelliSense \\
Data Grid        & \code{DataGrid.tsx}             & Sortable, paginated result display \\
Schema Explorer  & \code{SchemaExplorer.tsx}        & Tree view of tables, columns, indexes \\
ER Diagram       & \code{ERDiagram.tsx}             & ReactFlow entity-relationship graph \\
Query Builder    & \code{VisualQueryBuilder.tsx}    & GUI-driven SQL construction \\
Plan Visualiser  & \code{QueryPlanVisualizer.tsx}   & Tree view of EXPLAIN output \\
Monitoring       & \code{Monitoring.tsx}            & Buffer pool \& query metrics charts \\
Query History    & \code{QueryHistory.tsx}          & Searchable past-query list \\
Command Palette  & \code{CommandPalette.tsx}        & \texttt{Ctrl+K} fuzzy action launcher \\
\bottomrule
\end{tabular}
\end{table}

\subsection{Network Layer --- HTTP Server}

The C++ backend runs a \emph{zero-dependency} HTTP server built on raw BSD sockets (\code{server.h}).  It supports:
\begin{itemize}[nosep]
  \item Multithreaded request handling (one \code{std::thread} per connection).
  \item CORS headers for cross-origin browser requests.
  \item Routing to nine REST endpoints (see \cref{ch:web-interface}).
\end{itemize}

\subsection{Query Processing Layer}

\begin{figure}[H]
\centering
\begin{tikzpicture}[node distance=0.6cm and 1.2cm]
  \node[component, minimum width=2.2cm] (sql) {SQL String};
  \node[component, minimum width=2.2cm, right=of sql] (lex) {Tokeniser};
  \node[component, minimum width=2.2cm, right=of lex] (tok) {Tokens};
  \node[component, minimum width=2.2cm, right=of tok] (par) {Parser};
  \node[component, minimum width=2.2cm, right=of par] (ast) {AST};
  \node[component, minimum width=2.2cm, right=of ast] (exe) {Executor};

  \draw[arr] (sql) -- (lex);
  \draw[arr] (lex) -- (tok);
  \draw[arr] (tok) -- (par);
  \draw[arr] (par) -- (ast);
  \draw[arr] (ast) -- (exe);
\end{tikzpicture}
\caption{Query processing pipeline.}
\label{fig:query-pipeline}
\end{figure}

\begin{enumerate}[nosep]
  \item \textbf{Tokeniser} (\code{tokenizer.h}) --- scans the SQL string character-by-character, emitting a stream of typed tokens (keywords, identifiers, literals, operators).
  \item \textbf{Parser} (\code{parser.h}) --- a hand-written recursive-descent parser that consumes tokens and builds an Abstract Syntax Tree (AST).
  \item \textbf{Executor} (\code{executor.h}) --- walks the AST object and calls into the catalog and storage layers to produce a \code{QueryResult}.
\end{enumerate}

\subsection{Catalog Layer}

The \code{Catalog} class (\code{catalog.h}) is the system's metadata repository.  It:
\begin{itemize}[nosep]
  \item Stores \code{TableSchema} objects (column names, types, constraints).
  \item Manages \code{IndexMeta} records.
  \item Persists everything to \code{catalog.json} using the \code{nlohmann/json} library.
  \item Loads the catalog on startup and saves after each DDL operation.
\end{itemize}

\subsection{Storage Layer}

\begin{figure}[H]
\centering
\begin{tikzpicture}[node distance=1cm and 1.5cm]
  \node[component] (tbl) {Table};
  \node[component, below left=of tbl] (bp) {Buffer Pool};
  \node[component, below right=of tbl] (dm) {Disk Manager};
  \node[page, below=of bp] (pg) {Page\\(4 KB)};
  \node[storage, below=of dm] (file) {.nxdb file};

  \draw[arr] (tbl) -- (bp);
  \draw[arr] (tbl) -- (dm);
  \draw[arr] (bp) -- (pg);
  \draw[arr] (dm) -- (file);
  \draw[dasharr] (bp) -- (dm) node[midway, above, font=\scriptsize] {evict/fetch};
\end{tikzpicture}
\caption{Storage layer components.}
\label{fig:storage-layer}
\end{figure}

\begin{itemize}[nosep]
  \item \textbf{Page} (\code{page.h}) --- the 4\,KB unit of I/O; uses a slotted-page layout.
  \item \textbf{Disk Manager} (\code{disk\_manager.h}) --- reads/writes pages from \code{.nxdb} files.
  \item \textbf{Buffer Pool} (\code{buffer\_pool.h}) --- an LRU cache that keeps frequently-used pages in memory.
  \item \textbf{Table} (\code{table.h}) --- serialises/deserialises rows and orchestrates page-level operations.
\end{itemize}

\section{Data Flow Example}

\begin{deepdivebox}[Tracing \texttt{SELECT * FROM customers WHERE city = 'NYC'}]
\begin{enumerate}[nosep]
  \item The frontend sends \code{\{"sql": "SELECT ..."\}} via \code{POST /api/query}.
  \item \code{HttpServer::handleQuery} extracts the SQL string and calls \code{Executor::execute}.
  \item The \textbf{Tokeniser} produces tokens: \code{SELECT}, \code{STAR}, \code{FROM}, \code{IDENTIFIER(customers)}, \code{WHERE}, \code{IDENTIFIER(city)}, \code{EQUALS}, \code{STRING\_LITERAL('NYC')}.
  \item The \textbf{Parser} builds a \code{SelectStatement} AST with a \code{BinaryOp(EQ)} WHERE clause.
  \item The \textbf{Executor} looks up ''customers'' in the Catalog, gets the \code{Table} object.
  \item \code{Table::scanAll()} iterates over pages via the Buffer Pool.
  \item For each page, the Buffer Pool checks its LRU cache; on a miss, it calls \code{DiskManager::readPage} to load from the \code{.nxdb} file.
  \item  Each tuple is deserialized and tested against the WHERE predicate.
  \item Matching rows are collected into a \code{QueryResult} and serialised to JSON.
  \item The HTTP server sends the JSON response back to the React frontend.
\end{enumerate}
\end{deepdivebox}

\section{Namespace and Code Organisation}

All C++ code lives in the \code{shift\_elite} namespace.  The directory structure mirrors the architectural layers:

\begin{lstlisting}[style=pseudo, numbers=none]
backend/src/
  types/types.h          -- Value, DataType, TableSchema, constants
  storage/
    page.h               -- Page class (4 KB slotted page)
    disk_manager.h       -- File I/O for .nxdb files
    buffer_pool.h        -- LRU buffer pool
    table.h              -- Row serialisation, insert/scan/delete
  catalog/catalog.h      -- Schema and index metadata
  sql/
    tokenizer.h          -- SQL lexer
    ast.h                -- AST node hierarchy
    parser.h             -- Recursive-descent parser
    executor.h           -- Statement execution engine
  server/server.h        -- HTTP server + REST API routing
  main.cpp               -- Entry point, sample data seeding
\end{lstlisting}

\begin{interviewbox}[Interview Corner]
\textbf{Q:} Describe the architecture of a typical relational DBMS.  Which layers does \shiftdb{} implement?\\[4pt]
\textbf{A:} A textbook DBMS has a \emph{query processor} (parser + optimiser + executor), a \emph{storage engine} (buffer pool + disk manager), a \emph{catalog}, and a \emph{transaction manager}.  \shiftdb{} implements all of these except a full-blown optimiser (it uses a simple cost model in \code{EXPLAIN}).  It adds a presentation layer (React UI) and a network layer (HTTP server) on top.

\vspace{6pt}
\textbf{Q:} Why is separation of layers important in database design?\\[4pt]
\textbf{A:} Layer separation provides (1) \emph{modularity} --- each component can evolve independently; (2) \emph{testability} --- layers can be unit-tested in isolation; (3) \emph{portability} --- the storage layer can be swapped (e.g., from file system to cloud storage) without touching the parser.
\end{interviewbox}
